\chapter{Вычислительная лаборатория серии}
\label{ch:computational-laboratory}

\section{Затянувшаяся лаборатория, а не только теоремы}
\label{sec:lab-role}

Монография задаёт закон~$g$, геометрию носителя и цепочку $M\to T$.
Проверка этой цепочки для читателя устроена как \emph{затянувшаяся физическая лаборатория},
в которой вместо осциллографа --- симулятор решётки и алгебраические сверки,
а вместо лабораторного журнала --- явные \textbf{записи опытов} в конце каждого тома.

Три роли не смешиваются:
\begin{description}[style=nextline,leftmargin=3.2cm]
\item[\textbf{Теория (MODEL)}]
  что \emph{должно} следовать из посылок; без дат сессий и без статуса кода.
\item[\textbf{Аппарат}]
  дискретная реализация шага~$g$, приборы осреднения, сценарии начальных условий
  (в репозитории проекта --- пакет симуляции и сценарии запуска).
\item[\textbf{Журнал прогонов}]
  серия вычислительных экспериментов: полная постановка, наблюдаемые,
  анализ по \cref{ch:measurement-processing}, вывод и статус (закрыт / открыт / провал).
\end{description}

\begin{remark}[Не подмена натурного опыта]
Вычислительный прогон проверяет \emph{согласованность} заявленного $g$ и моста $M\to T$
на заданной постановке.
Сравнение с PDG/CODATA остаётся на~$\Tlayer$ и требует метрологии эталона
(\cref{ch:measurement-processing}), а не «красивого скриншота поля''.
\end{remark}

\section{Бремя доказательства}
\label{sec:lab-burden-of-proof}

Утверждения серии сильные, потому что вывод жёсткий: \textbf{полное квантование на~$M$},
две скорости $c_0$ и $c$, лестница $E_0$, перепись $M=97$ --- не метафора и не «авторы решили, что ок''.
Публикация сама по себе никому ничего не доказывает: для каждого класса утверждений мы обязаны дать
хотя бы одно из:
\begin{enumerate}[label=\textbf{B\arabic*:},leftmargin=*]
\item \textbf{доказательство} из явных посылок (теорема с гиперссылкой на аксиому);
\item \textbf{запись опыта} (серия CE-*) с наблюдаемой, порогом PASS/FAIL и командой воспроизведения;
\item \textbf{явный статус «открыто''} с критерием decisive-провала (что именно сломает модель).
\end{enumerate}
Сравнение с CODATA/PDG без строки Verify и без \cref{ch:measurement-processing} --- недопустимая
замена опыта.
Рецензент не обязан «догадываться'', что имелось в виду: если в тексте нет CE-записи или реестра,
утверждение считается \emph{не проверенным в рамках серии}, даже если формула выглядит правдоподобно.

\section{Единый протокол записи опыта}
\label{sec:lab-protocol}

Каждая запись в журнале тома следует одному шаблону
(макросы \texttt{labrecord} в преамбуле книги):

\begin{labrecord}{Шаблон (пустая форма)}{lab:template}
\labgoal{Какой вопрос MODEL ставит на $\Tlayer$ или на sim-прокси.}
\labhypothesis{Численное или качественное предсказание \emph{до} прогона.}
\labmodel{Пункты MODEL (§, теорема), которые опыт должен поддержать или опровергнуть.}
\labsetup{Носитель, размер, границы, начальное состояние, число тактов, прибор чтения.}
\labobservables{Что именно извлекается из поля (норма, спектр, ось, время жизни прокси, \ldots).}
\labcriterion{PASS/FAIL: пороги, допуски, $E_n$ к эталону, ансамбль + Стьюдент.}
\labanalysis{Тип~A/B, перенос $u$, отделение систематики модели от табличного $u$.}
\labconclusion{Подтверждено / частично / не проверено / опровергнуто; что переносится в следующий опыт.}
\labreproduce{Идентификатор в реестре автоматических сверок; параметры сценария для повтора.}
\end{labrecord}

Серии опытов (\texttt{labseries}) объединяют прогоны с общим аппаратом
(один и тот же кипящий вакуум, одна цепочка аннигиляции, одна лестница~$\alpha$).

\section{Связь с метрологией и критическими экспериментами}
\label{sec:lab-metrology-link}

\begin{enumerate}[label=\textbf{L\arabic*:},leftmargin=*]
\item \textbf{Постановка} фиксирует, что является наблюдаемой, а что --- артефактом дискретизации.
\item \textbf{Анализ} --- по \cref{ch:measurement-processing}: прямые/косвенные величины, GUM, Стьюдент для ансамбля.
\item \textbf{Критические эксперименты} в конце тома --- не буллет-лист, а развёрнутый журнал
  (\CiteConstruction{ch:critical-construction}; журнал этажей --- том \emph{floors}; аналоги в прочих томах).
\item Опыт со статусом \emph{открыт} остаётся в журнале с явным критерием decisive-провала модели.
\end{enumerate}

\section{Классы вычислительных опытов}
\label{sec:lab-classes}

\begin{description}[style=nextline,leftmargin=3.4cm]
\item[\textbf{Алгебра / замыкание}]
  SI-строка, лестница $\alpha$, тождества масс без поля на решётке.
  Критерий --- внутренняя согласованность и $E_n$ к CODATA при заявленных $u$.
\item[\textbf{Макро-поле на карте}]
  дисперсия, фронт, скорость $c$ после $\mathcal{B}$ (\CiteConstruction{ch:macro}).
\item[\textbf{Составной узел на этаже}]
  два- и $n$-телесные сценарии, оболочки, аннигиляция, прокси диполя и времени жизни.
\item[\textbf{Отрицательные / контроль}]
  сценарий, который \emph{не} должен давать эффект (контроль pp вместо pm и т.\,п.).
\end{description}

Журнал пополняется по мере закрытия UoW; закрытая запись не переписывает MODEL,
а фиксирует, какой класс проверок на данном этапе считается выполненным.
